%!TEX root = ../double_trouble_supp.tex

In this section, we provide additional details and proofs for the results presented in Section 4.
We state and prove a slight generalization of Theorem 1: in our proof the number of populations can be two (as in the statement presented in the paper) or greater: $\numpops > 1$. 
The two-population case of Theorem 1 follows as a corollary.

Before stating our result, we extend our notation to encompass any strictly positive integer number of populations $\numpops \in \{1,2,\dots\}$.
We let $\samplesize{\popidx}$ denote the number of observed pilot samples in population $\popidx$ with $1 \le \popidx \le P$.
We collect all the pilot sizes in the vector $\vecsamplesize := (\samplesize{1}, \ldots,\samplesize{\numpops})$.
We collect the observations from the $\unitidx$-th sample in the $\popidx$-th population in $\genericmeasurevariantcount$ as in Equation (2.1).
We write $\measurevariantcount{\popidx}{1:\samplesize{\popidx}}$ for the pilot samples in population $\popidx$.
And we write $X_{1:\numpops, \bm{1}:\vecsamplesize}$ for the full collection of pilot samples across all $\numpops$ populations.

We collect the frequencies associated with the $\variantidx$-th variance across the $\numpops$-populations as $\genericvecvariantfreq = (\variantfreq{1}{\variantidx}, \ldots, \variantfreq{\numpops}{\variantidx})$.
The vector-valued measure $\randommeas :=  \sum_{\variantidx=1}^{\infty} \genericvecvariantfreq \dirac{\genericvariantlabel}$ matches variant frequency
vectors with their labels. We write $\randommeas \sim \PPP(\levync)$ to indicate that $\randommeas$
is generated by drawing $\{\genericvecvariantfreq\}_{\variantidx=1}^{\infty}$ from a Poisson
point process (PPP) with rate measure $\levync(\de\vecvariantfreq{})$ and drawing the $\genericvariantlabel$
i.i.d.\ uniform on $[0,1]$. We model $\genericvariantcount \sim \bernoullirv{\genericvariantfreq}$, i.i.d.\ across $\unitidx$
and independent across $\popidx$ and $\variantidx$ \revision{conditioning on $\genericvariantfreq$}; when the $\genericvariantcount$ are drawn this way, we say that
$X_{1:\numpops, \bm{1}:\vecsamplesize}$
is drawn according to a \emph{multiple-population Bernoulli process} ($\mBeP$) with measure parameter $\randommeas$ and count-vector parameter $\vecsamplesize$.
The $\numpops$-population generalization of Equation (3.1) is then
\begin{align}
	\begin{split}
		\randommeas & \sim \PPP(\nu) \quad \textrm{ (Prior) } \\
		X_{1:\numpops, \bm{1}:\vecsamplesize} \mid \randommeas & \sim \mBeP(\randommeas, \vecsamplesize). \quad \textrm{ (Likelihood) }
		\label{eq:appx_generative_model}
	\end{split}
\end{align}
In particular, the notation and development before are the $\numpops=2$ special case of what is written above.
Observe that Desideratum (A) and Desideratum (B) do not specifically reference $\numpops=2$ and can be
applied to the general $\numpops > 0$ case as well.

To prove our main result, we will find it useful to show that each of Desideratum (A) and Desideratum (B) is equivalent to
an integral constraint on the L\'{e}vy  rate measure $\levync$ of the prior distribution $\randommeas$. To that end, we state and prove the following lemma.
\begin{myLemma}[Integral requirements]
	Take a number of populations $\numpops$ with $1 < \numpops < \infty$.
	Take data generated according to the model in \cref{eq:appx_generative_model}. Then Desideratum (A) holds if and only if the rate measure $\levync$ satisfies  
	\begin{equation}
		\forall \popidx \in \{1,\ldots,\numpops\}, \quad \int_{[0,1]^\numpops} \variantfreqperpop{\popidx} \levync(\de\vecvariantfreq{})<\infty.  \label{eq:finite_first_moment}
	\end{equation} 
	And Desideratum (B) holds if and only if the rate measure $\levync$ satisfies 
	\begin{equation}
		\forall \popidx \in \{1,\ldots,\numpops\}, \quad \int_{(0,1]}  \levync_{-\popidx} (\de \variantfreqperpop{\popidx})=\infty. \label{eq:infinite_mass}
	\end{equation}
		\internal{Here $\levync_{-\popidx} (\de \variantfreqperpop{\popidx}) := \int_{\vecvariantfreq{-\popidx}\in [0,1]^{\numpops-1}} \levync(\de\vecvariantfreq{})$, where $\vecvariantfreq{-\popidx} = (\variantfreqperpop{1},\ldots,\variantfreqperpop{\popidx-1}, \variantfreqperpop{\popidx+1}, \ldots,\theta_\numpops)$ denotes a vector containing all population indices other than $\popidx$.}
    \label{lemma:integral_requirements}
\end{myLemma}
\begin{proof}
	By construction, a point drawn from the Poisson point process with rate measure $\levync$ with support on $[0,1]^\numpops$ is of the form 
	$\vecvariantfreq{} = (\variantfreqperpop{1}, \ldots,\theta_\numpops) \in [0,1]^\numpops$. 
	By standard properties of multivariate Poisson processes, the $\popidx$-th coordinate of the multivariate Poisson process is a univariate Poisson point process with rate measure $\levync_{-\popidx} (\de \variantfreqperpop{\popidx}) := \int \levync(\de\vecvariantfreq{-\popidx})$ \citep[Mapping theorem,][ch. 2.3]{kingman1992poisson}.

	In our setting, then, an individual $\genericmeasurevariantcount$ in population $\popidx$ is drawn from a Bernoulli process, where the parameter encoding the variants with their frequency follows a Poisson point process with rate measure $\levync_{-\popidx}$. Following the same argument in \citet[][Section 2.3]{Broderick2018}, $\genericmeasurevariantcount$ shows finitely many variants (Desideratum (A)) with probability one if and only if the underlying rate measure has finite expectation on $[0,1]$. In our case then, population $\popidx$'s frequencies are obtained from the marginal rate $\levync_{-\popidx}$; thus individuals exhibit finitely many variants if and only if $\int\theta_\popidx \levync_{-\popidx}(\de \variantfreqperpop{\popidx}) < \infty$.
	By substituting the definition of $\levync_{-\popidx}(\de \variantfreqperpop{\popidx})$, $\int\theta_\popidx \levync_{-\popidx}(\de \variantfreqperpop{\popidx}) < \infty$ can be rewritten as \cref{eq:finite_first_moment}.
	 

	Desideratum (B) is satisfied if and only if for each population there are infinitely many variants that can be discovered. Again by \citet[][Section 2.3]{Broderick2018} population $\popidx$ has infinitely many variants to be seen if and only if $\levync_{-\popidx} (\de \variantfreqperpop{\popidx})$ has infinite mass on $(0,1]$, which is the requirement stated in \cref{eq:infinite_mass}.  
\end{proof}

\begin{myRemark}[On Desideratum (B)]
	\label{remark:inf_mass}
	A traditional requirement of one-population Bayesian nonparametric feature models is that the rate measure has infinite mass on $(0,1]$ \citep{Gnedin2007,Broderick2012}.
	Desideratum (B) provides an extension of this requirement to multiple populations. If we 
	think of Desideratum (B) as an extension from a single dimensional $\variantfreq{}{}$ to a multidimensional $ \vecvariantfreq{}$, 
	there are a number of other possible extensions that might a priori seem reasonable.
	For instance, we might instead require infinite mass on $(0,1]^\numpops$ (i.e., all vectors in the unit cube with strictly positive entries) or
	infinite mass on $[0,1]^\numpops\setminus{\bm{0}}$ where $\bm{0}$ is a $\numpops$-long 0-vector (i.e., the unit cube without the origin).
	
	We next observe that these alternative extensions are not equivalent to Desideratum (B).
	To show how these are not equivalent, we next provide simple counterexamples. To help us do so, we will use $\numpops$ (unidimensional) measures $\mu_{1}(\de\variantfreqperpop{1}),$ $\dots ,\mu_{\numpops}(\de\variantfreqperpop{\numpops})$. We assume that for $\popidx=1,\dots,\numpops$, we have $\int_{(0,1]} \mu_{\popidx}(\de\variantfreqperpop{\popidx})=\infty$ and $\int_{(0,1]}\variantfreqperpop{\popidx} \mu_{\popidx}(\de\variantfreqperpop{\popidx})<\infty$. We denote the delta measure at $0$ as $\delta_0(\cdot)$. 
	
%	First, we construct a measure that has 0 mass on $(0,1]^\numpops$ but satisfies Desideratum (B). To construct our measure, we put mass only on the axes; namely, we consider the measure $\sum_{\popidx=1}^\numpops\mu_{\popidx}(\de\variantfreqperpop{\popidx})\prod_{k: k\ne \popidx}\delta_{0}(\de\variantfreqperpop{k})$. This measure can be interpreted as the law describing the behavior of $\numpops$ populations that do not share variants; each population separately draws its variant frequencies from $\mu_{\popidx}$. Since by assumption $\mu_{\popidx}$ satisfies the single-population requirements, the full measure satisfies Desideratum (B). But the measure has 0 mass on $(0,1]^\numpops$ because its mass lies entirely along the ``axes''.

	As a first counterexample, we construct a measure that has 0 mass on $(0,1]^\numpops$ but satisfies Desideratum (B). 
	To construct our measure, we put mass only on the coordinate axes of the $\numpops$-dimensional unit cube $[0,1]^\numpops$; 
	namely, we consider the measure $\sum_{\popidx=1}^\numpops\mu_{\popidx}(\de\variantfreqperpop{\popidx})\prod_{k: k\ne \popidx}\delta_{0}(\de\variantfreqperpop{k})$. 
	Here, the coordinate axes refer to the sets $\{(\variantfreqperpop{1}, 0, \ldots, 0) : \variantfreqperpop{1} \in [0,1]\}$, $\{(0, \variantfreqperpop{2}, 0, \ldots, 0) : \variantfreqperpop{2} \in [0,1], \dots \}$, where exactly one coordinate is non-zero. 
	This measure can be interpreted as the law describing the behavior of $\numpops$ populations that do not share variants; each population separately draws its variant frequencies from $\mu_{\popidx}$, and variants appear in exactly one population since they have positive frequency along only one coordinate axis. 
	Since by assumption $\mu_{\popidx}$ satisfies the single-population requirements, the full measure satisfies Desideratum (B). 
	But the measure has 0 mass on $(0,1]^\numpops$ because its mass lies entirely along these coordinate axes, never in the interior of the unit cube where multiple coordinates would be simultaneously positive (which would correspond to shared variants).

%	As a second counterexmaple, we construct a measure that has infinite mass on $[0,1]^\numpops\setminus{\bm{0}}$ but does not satisfy Desideratum (B). In this case, we put mass on exactly one axis. Then all but one population exhibit zero variants. In particular, we consider the measure $\mu_{1}(\de\variantfreqperpop{1})\prod_{k=2}^\numpops\delta_0(\de\variantfreqperpop{k})$. This measure has infinite mass on $[0,1]^\numpops\setminus{\bm{0}}$. But when we sample from, e.g., population 2, we will not see any variants since all variants have frequency 0. Thus this measure does not satisfy Desideratum (B). 
	As a second counterexample, we construct a measure that has infinite mass on $[0,1]^\numpops\setminus{\bm{0}}$ but does not satisfy Desideratum (B). 
	In this case, we put mass on exactly one coordinate axis, specifically the first axis. 
	Then all but one population exhibit zero variants. 
	In particular, we consider the measure $\mu_{1}(\de\variantfreqperpop{1})\prod_{k=2}^\numpops\delta_0(\de\variantfreqperpop{k})$, which places all mass on the set $\{(\variantfreqperpop{1}, 0, 0, \ldots, 0) : \variantfreqperpop{1} \in [0,1]\}$. 
	This measure has infinite mass on $[0,1]^\numpops\setminus{\bm{0}}$ since the first coordinate axis (excluding the origin) is a subset of this set and $\mu_1$ has infinite mass on $(0,1]$ by assumption. 
	However, this measure violates Desideratum (B) because when we sample from any population $\popidx \geq 2$, we will not see any variants since all variants have frequency $\variantfreqperpop{\popidx} = 0$ in those populations. 
	In other words, only population 1 exhibits variants while populations 2 through $\numpops$ are effectively "empty," which contradicts the requirement that each population should have infinitely many discoverable variants.


\end{myRemark}


We now extend the statement of \cref{prop:incomp} to the general setting for an arbitrary number of populations $\numpops$. 
To do so, we first restate~Condition (C), now for a generic number of populations.
\begin{conditionp}{C'} \label{req:appx_product_measure}
	The rate measure for the frequency-generating PPP in $\numpops$ populations factorizes across populations:
	$\levync(\de\btheta) = \prod_{\popidx=1}^{\numpops} \levync_{\popidx}(\de\variantfreqperpop{\popidx})$.  
\end{conditionp}
Notice that Condition (C) is a special case of the condition above, when $\numpops = 2$.

We can now state the result which generalizes Theorem 1.

\begin{myTheorem} \label{prop:incomp_general}
	Take the general $\numpops$ population model of \cref{eq:appx_generative_model}. If the Poisson point process
	rate measure in the prior factorizes across populations according to \cref{req:appx_product_measure},
	the model cannot simultaneously generate samples with finitely many variants (Desideratum (A))
	and guarantee that there are always more variants to discover (Desideratum (B)).
\end{myTheorem}

\begin{proof}

To prove the result, it is enough to assume any two of Desideratum (A), Desideratum (B), and \cref{req:appx_product_measure} and then show that the remaining one cannot hold. 

We show that if Desideratum (B) and \cref{req:appx_product_measure} hold, then Desideratum (A) cannot hold:
\[
	\text{If }Desideratum (B) \text{ and } \Cref{req:appx_product_measure} \text{ hold, } Desideratum (A) \text{ does not hold}.
\]
By \cref{req:appx_product_measure} we have that $\levync$ is a ``product measure'', i.e. $\levync(\de\vecvariantfreq{})=\prod_{\popidx} \levync_{\popidx}(\de\variantfreqperpop{\popidx})$ and in turn we have $\levync_{-\popidx} (\de \variantfreqperpop{\popidx}) = \levync_{\popidx}(\de\variantfreqperpop{\popidx}) \prod_{k: k\ne \popidx} \int_{[0,1]}\levync_{k}(\de\variantfreqperpop{k})$. By Desideratum (B) and its equivalent~\cref{eq:infinite_mass}, for all $\popidx$, we have $\int_{(0,1]} \levync_{-\popidx} (\de \variantfreqperpop{\popidx})=\infty$ thus either:
\begin{itemize}
	\item[(i)] $\int_{(0,1]}\levync_{\popidx}(\de\variantfreqperpop{\popidx})=\infty$ and all other $k\ne \popidx$, $\int_{[0,1]}\levync_{k}(\de\variantfreqperpop{k})>0$, or 
	\item[(ii)] $0<\int_{(0,1]}\levync_{\popidx}(\de\variantfreqperpop{\popidx})<\infty$ and there exists at least one $k\ne \popidx$ for which $\int_{[0,1]}\levync_{k}(\de\variantfreqperpop{k})=\infty$, \revision{with $h\ne \popidx$ and $h \ne k$ we have $\int_{[0,1]}\levync_{h}(\de\variantfreqperpop{h})>0$}.
\end{itemize}
We now show that in both cases (i), (ii),  Desideratum (A) does not hold.

In case (i), for $k\ne \popidx$:
\[
	\int_{[0,1]^{\numpops}}\variantfreqperpop{k}\levync(\de\vecvariantfreq{}) = \int_{[0,1]}\variantfreqperpop{k} \levync_{k}(\de\variantfreqperpop{k}) \int_{[0,1]} \levync_{\popidx}(\de\variantfreqperpop{\popidx})\prod_{i: i\ne k, \popidx} \int_{[0,1]}\levync_{i}(\de\variantfreqperpop{i})=\infty.
\]
Since we have $\prod_{i: i\ne k, \popidx} \int_{[0,1]}\levync_{i}(\de\variantfreqperpop{i})>0$ and $\int_{[0,1]}\levync_{\popidx}(\de\variantfreqperpop{\popidx})=\infty$, it remains to check $\int_{[0,1]} \variantfreqperpop{k} \levync_{k}(\de\variantfreqperpop{k})>0$. 
Recall that we assumed~\cref{eq:infinite_mass} is satisfied, we must have the ``projected'' measure $\levync_{k}(\de\variantfreqperpop{k}) \int_{[0,1]} \levync_{\popidx}(\de\variantfreqperpop{\popidx})\prod_{i: i\ne k, \popidx} \int_{[0,1]}\levync_{i}(\de\variantfreqperpop{i})$ to have mass on $(0,1]$, which implies $\levync_{k}(\de\variantfreqperpop{k})$ must have mass on $(0,1]$. Then we must also have $\int_{[0,1]}\variantfreqperpop{k} \levync_{k}(\de\variantfreqperpop{k})>0$, hence implying that Desideratum (A) does not hold.

In case (ii), by direct factorization of the rate measure it holds that
\[
	\int_{[0,1]}\variantfreqperpop{\popidx}\levync(\de\vecvariantfreq{}) = \int_{(0,1]}\variantfreqperpop{\popidx} \levync_{\popidx}(\de\variantfreqperpop{\popidx}) \prod_{k: k\ne \popidx} \int_{[0,1]}\levync_{k}(\de\variantfreqperpop{k})=\infty.
 \] 
Hence, also in this case, Desideratum (A) does not hold. We have then showed that if Desideratum (B) and \cref{req:appx_product_measure} hold, then Desideratum (A) cannot hold.

\end{proof}

\begin{proof}[Proof of Theorem 1]
	Theorem 1 is a special case of \cref{prop:incomp_general} when $\numpops = 2$. 
\end{proof}

